% Table 6 of paper_bc_peerjcs.tex, exported for PeerJ upload by
% scripts/pubs/export_peerj_upload.py. Do not edit: edit the manuscript and re-export.
\documentclass[10pt]{article}
\usepackage[letterpaper,margin=25mm]{geometry}
\usepackage{amsmath,amssymb}
\usepackage{booktabs}
\usepackage{tabularx}
\usepackage{array}
\usepackage{multirow}
\usepackage{enumitem}
\usepackage[hidelinks]{hyperref}
\usepackage[authoryear,round]{natbib}
\begin{document}
\begin{table}[t]
\centering
\caption{Core metric families in model-agnostic XAI evaluation: what
each captures well and what it leaves unresolved.}
\footnotesize
\begin{tabularx}{\linewidth}{l X X}
\toprule
\textbf{Metric family} & \textbf{Captures well} & \textbf{Leaves unresolved} \\
\midrule
Fidelity / faithfulness & Whether explanations track model behavior under interventions or masking. & Whether the explanation is meaningful, plausible, or useful outside the proxy setup. \\
Stability / robustness & Whether explanations remain consistent across perturbations or repeated runs. & Whether stable explanations are actually correct or semantically adequate. \\
Sparsity / parsimony & Whether explanations are concise and low in feature load. & Whether concision hides necessary nuance or context. \\
Benchmark-grounded correctness & Whether explanations recover synthetic or domain-defined ground truth. & Whether the ground truth is realistic, complete, or transferable. \\
Human expert review & Whether explanations align with domain expectations and expert reasoning. & Whether those judgments are scalable, standardized, and reproducible. \\
User-centered evaluation & Whether explanations improve understanding, decisions, or trust calibration in practice. & Whether findings generalize across tasks, populations, and study designs. \\
LLM-based semantic evaluation & Whether rubric-based semantic judgments can be produced quickly and at scale. & Whether those judgments agree with humans and avoid proxy-judge bias. \\
\bottomrule
\end{tabularx}

\end{table}
\end{document}
